% <!--
% 文件：main.tex
% 核心功能：DaE 分支论文 v2 LaTeX，empirical method paper，针对 advisory 子集（30 personas, 5 sub-scenarios）。
% 输入：v1.0 实验数据 advisory 子集（subset_advisory.json）、refs.bib、AKM 母论文交叉引用。
% 输出：可投 arXiv (cs.HC) 的 8-10 页论文。
% -->
\documentclass[11pt]{article}

\usepackage[margin=1in]{geometry}
\usepackage[T1]{fontenc}
\usepackage[utf8]{inputenc}
\usepackage{lmodern}
\usepackage{setspace}
\usepackage{enumitem}
\usepackage{booktabs}
\usepackage[numbers]{natbib}
\usepackage[hidelinks]{hyperref}

\setstretch{1.12}
\setlength{\parskip}{0.5em}
\setlength{\parindent}{0pt}

\title{Cold-Start Profile Recovery in LLM Advisory: A Cross-Family Three-Judge Benchmark on 30 Synthetic Advisory Personas}
\author{Weishi Shao\thanks{Independent researcher; corresponding repository: \url{https://github.com/sirsws/akm} under \texttt{branches/dae}.}}
\date{May 2026}

\begin{document}

\maketitle

\begin{abstract}
Most AI advisory systems begin generating advice before they have constructed any explicit model of the operator. The result is alignment debt: users repeatedly restate goals and constraints, and the system defaults to generic guidance. We study \textbf{Dialogue-as-Elicitation (DaE)}, an instance of the broader Active Knowledge Modeling (AKM) paradigm \citep{shao2026akm}, which converts the cold-start moment into an explicit elicitation pass that yields a structured, reusable user profile. We evaluate DaE on 30 synthetic personas covering five advisory sub-domains (career, learning, investment, health/lifestyle, relationships) under four pre-task context conditions (\textit{no\_profile}, \textit{unstructured\_notes}, \textit{akm\_profile}, \textit{akm\_elicited}) and three independent cross-family LLM judges (DeepSeek-V4-pro, Gemini-3-Flash, Grok-4.3) blindly scoring the same anonymized 4-answer set per persona. Pooled total scores (out of 30) are 10.80, 21.06, 29.43, 29.10 respectively (n${=}$90 per condition). Inter-judge reliability on the advisory subset is high in aggregate (Krippendorff $\alpha = 0.950$ on totals; pairwise Spearman $\rho = 1.0$ on the four-condition ranking across all judge pairs), with one per-dimension caveat (risk control $\alpha = 0.78$, slightly below the conventional $0.80$ line; discussed in Section~\ref{sec:method}). The DaE-elicited profile recovers $98.9\%$ of the hand-curated upper bound on this saturated 1--5 rubric. We read this as a reproducible empirical signal that, under controlled synthetic personas and three cross-family LLM judges, an automatic single-pass cold-start elicitation can match a curated profile on six advisory-quality dimensions; we are explicit about the synthetic-only ceiling and the limitations of LLM-as-judge on saturated rubrics.
\end{abstract}

\textbf{Keywords:} cold-start profile recovery, LLM advisory benchmark, cross-family LLM-as-judge, inter-rater reliability, Krippendorff alpha, pre-task user-profile conditioning, active knowledge modeling.

\section{Introduction}

In advisory settings --- career, learning, finance, health, and relationships --- the bottleneck is rarely raw model capability. It is alignment transfer: the user's goals, hard constraints, and decision style are tacit and expensive to communicate through ad-hoc prompting. Recent work names this hidden labor \emph{alignment debt} \citep{oyemike2025alignment}.

\textbf{Dialogue-as-Elicitation (DaE)} treats the first interaction differently. Rather than answering immediately, the system runs a short structured elicitation pass and produces a reusable structured profile that can be loaded into any subsequent advisory turn. The DaE artifact contains the same six fields used elsewhere in the AKM paradigm \citep{shao2026akm}: goals, hard constraints, soft preferences, available assets, risk boundaries, decision style.

The empirical question is sharp: \emph{does the cold-start DaE pass actually produce a profile good enough to drive downstream advice quality?} An earlier social-science framing of DaE laid out the mechanism design \citep{shao2025dae}; this paper provides the missing empirical answer. The contribution is twofold:

\begin{enumerate}[leftmargin=1.5em]
\item A reproducible cold-start evaluation across 30 synthetic personas spanning five advisory sub-domains, with four pre-task context conditions and three cross-family LLM judges.
\item A reliability-grounded result: $\alpha = 0.950$ across three judges on totals, $\rho = 1.0$ on the four-condition ranking for all judge pairs, and $98.9\%$ recovery of the hand-curated upper bound by the DaE-elicited profile on the saturated 1--5 rubric.
\end{enumerate}

We are explicit about the synthetic-only design (Section~\ref{sec:limits}) and what it cannot license. The bounded claim is: under controlled synthetic personas spanning advisory contexts and three cross-family LLM judges in near-perfect agreement, a single-pass DaE elicitation produces a profile that matches a hand-curated profile on six downstream advisory-quality dimensions.

\section{The DaE Mechanism}

\subsection{The Cold-Start Problem}

Advisory advice depends on operator state that is not in the prompt. Constraints (\emph{cannot quit job; budget under 3000 RMB; no late evening meetings}) and risk boundaries (\emph{no naked exits; no overexposed personal cases}) are typically unstated in a single user message. Most assistants either guess (and produce generic advice) or ask one disconnected clarifying question and proceed. Both paths leak alignment debt across sessions.

\subsection{The Elicitation Pass}

DaE treats the cold start as a single short pass: 5--7 targeted follow-up questions covering the six AKM fields, simulated user answers, then structured profile extraction. The output is a JSON profile that downstream advisory turns can consume without re-asking.

\subsection{Why It Should Work}

Three forces predict DaE will be useful even on a cold start:

\begin{enumerate}[leftmargin=1.5em]
\item \textbf{Mechanism design.} A short structured questionnaire externalizes sticky information faster than ad-hoc prompting, which is the core insight of preference elicitation \citep{foschini2025preference} generalized from local preference queries to a complete operator-model construction protocol.
\item \textbf{Asset reuse.} The structured profile is portable across sessions, models, and tools. Effort spent on the elicitation pass is amortized against all downstream use.
\item \textbf{Constraint addressability.} A structured profile makes each constraint individually addressable in the downstream prompt, which existing personalization benchmarks suggest improves task fit \citep{salemi2023lamp,white2023prompt,li2016persona}.
\end{enumerate}

The empirical question is whether these forces survive the \emph{cold start}, when the system has never seen the user before and must build the profile from a single dialogue.

\section{Method}
\label{sec:method}

We use the AKM Benchmark Toolkit v1.0 \citep{shao2026akm}; the present paper analyzes the 30-persona advisory subset.

\subsection{Conditions}

Four pre-task context conditions, holding generator and prompt template constant:

\begin{enumerate}[leftmargin=1.5em]
\item \textit{no\_profile}: task only.
\item \textit{unstructured\_notes}: task + 80--160 character first-person self-description.
\item \textit{akm\_profile}: task + hand-curated structured AKM profile.
\item \textit{akm\_elicited}: task + structured profile generated automatically by the DaE elicitation pass (5--7 simulated questions, simulated user answers grounded in the hidden persona, then profile extraction). The hand-curated profile is never seen by the elicitation chain.
\end{enumerate}

The contrast \textit{akm\_elicited} vs \textit{akm\_profile} is the core DaE test.

\subsection{Personas}

30 synthetic personas, six per sub-domain (Table~\ref{tab:dae_personas}). Each persona contains a hidden ground-truth state, a hand-aligned structured profile, a first-person fragment, and a downstream task drawn from a sub-scenario template (e.g.\ \emph{plan the next three months}, \emph{evaluate whether to take the new project}).

\begin{table}[h]
\centering
\caption{Advisory persona distribution (n${=}$30).}
\label{tab:dae_personas}
\begin{tabular}{lcl}
\toprule
Sub-domain & n & Sample tasks \\
\midrule
career           & 6 & 3-month plan, low-risk pivot, project acceptance, energy reallocation \\
learning         & 6 & 3-month learning plan, certification ROI, year uplift, work-study \\
investment       & 6 & 12-month allocation, in/out timing, education fund, housing, passive income \\
health/lifestyle & 6 & improvement plan, work-day routine, weekend, top-risk, exam decision \\
relationship     & 6 & family repair, parental boundary, intimacy, kids, workplace, friends \\
\bottomrule
\end{tabular}
\end{table}

\subsection{Generator and Judges}

Generator: \texttt{deepseek-v4-pro} (thinking disabled, temperature 0.2, max tokens 1500). Three cross-family judges: \texttt{deepseek-v4-pro}, \texttt{google/gemini-3-flash-preview}, \texttt{x-ai/grok-4.3}. Each persona's four answers are presented to all three judges in the \emph{same} deterministic anonymized order (seeded by \texttt{profile\_id}). Each judge scores each anonymous answer on six 1--5 dimensions (constraint adherence, risk control, specificity, actionability, personal fit, trade-off awareness) and identifies a best answer.

\subsection{Reliability Statistics}

Krippendorff $\alpha$ (interval, three raters) on per-(persona, condition) totals; per-dimension $\alpha$; pairwise Spearman $\rho$ on each judge's mean condition ranking. Standard for ordinal/interval LLM-as-judge designs.

\subsection{Reproducibility}

Full pipeline (personas, prompts, generator, three-judge runner, aggregator with $\alpha$ and $\rho$, advisory subset extractor) at \url{https://github.com/sirsws/akm} under \texttt{experiments/synthetic\_profile\_eval/}. Subset stats reproducible via:

\begin{verbatim}
python subset_stats.py --tag advisory
\end{verbatim}

\section{Results}
\label{sec:results}

\subsection{Pooled Scores on the Advisory Subset}

Table~\ref{tab:advisory_pooled} reports pooled means across the three judges and 30 advisory personas (n${=}$90 per condition).

\begin{table}[h]
\centering
\caption{Advisory subset, pooled means by condition (n${=}$90; max total ${=}30$).}
\label{tab:advisory_pooled}
\small
\begin{tabular}{lcccccccc}
\toprule
Condition & Total & Constraint & Risk & Specificity & Actionability & Fit & Trade-off & SD(total) \\
\midrule
no\_profile          & 10.80 & 1.43 & 2.44 & 1.94 & 2.26 & 1.33 & 1.40 & 4.17 \\
unstructured\_notes  & 21.06 & 3.69 & 4.09 & 3.41 & 3.86 & 3.41 & 2.66 & 3.98 \\
akm\_profile         & 29.43 & 4.94 & 4.96 & 4.96 & 4.94 & 4.90 & 4.84 & 1.32 \\
akm\_elicited        & 29.10 & 4.88 & 4.93 & 4.89 & 4.84 & 4.86 & 4.81 & 1.53 \\
\bottomrule
\end{tabular}
\end{table}

The DaE-elicited cold-start condition recovers $98.9\%$ of the hand-curated upper bound on totals ($29.10/29.43$), with the largest single-dimension gap at $0.13$ on \emph{trade-off awareness}. Both AKM conditions sit roughly $8.4$ points above the unstructured-notes baseline on a 30-point scale, and roughly $18.5$ points above the no-profile baseline.

\subsection{Inter-Judge Reliability}

Table~\ref{tab:dae_agreement} reports the reliability statistics on the advisory subset alone.

\begin{table}[h]
\centering
\caption{Inter-judge reliability on the advisory subset (3 cross-family judges, 30 personas).}
\label{tab:dae_agreement}
\begin{tabular}{lc}
\toprule
Statistic & Value \\
\midrule
Krippendorff $\alpha$ on totals (interval, 3 raters) & 0.950 \\
$\alpha$ constraint\_adherence & 0.942 \\
$\alpha$ risk\_control         & 0.779 \\
$\alpha$ specificity           & 0.909 \\
$\alpha$ actionability         & 0.856 \\
$\alpha$ personal\_fit         & 0.943 \\
$\alpha$ tradeoff\_awareness   & 0.922 \\
Spearman $\rho$ DeepSeek-V4-pro vs Gemini-3-Flash & 1.00 \\
Spearman $\rho$ DeepSeek-V4-pro vs Grok-4.3       & 1.00 \\
Spearman $\rho$ Gemini-3-Flash  vs Grok-4.3       & 1.00 \\
\bottomrule
\end{tabular}
\end{table}

\subsection{Best-Answer Votes}

Each judge identifies one best answer per persona; pooled across 90 selections (30 personas $\times$ 3 judges) on the advisory subset, \textit{akm\_profile} receives 43 votes, \textit{akm\_elicited} 28, \textit{unstructured\_notes} 1, \textit{no\_profile} 0; the remaining 18 selections were not parseable as a single answer label in one judge's structured output and are excluded from the count without affecting the dimension scores in Table~\ref{tab:advisory_pooled}.

\section{Discussion}

\subsection{Cold-Start Recovery Is Robust}

The DaE-elicited profile recovers $98.9\%$ of the curated upper bound on totals, and on no individual dimension does the gap exceed $0.13/5$. This is the central practical finding: a single short structured elicitation pass can produce, on a cold start, a profile that downstream advisory generation uses about as well as a hand-aligned profile would.

\subsection{Why Three Judges Matter}

Single-model LLM-as-judge designs are vulnerable to family-specific stylistic preferences. Our three judges are from three distinct model families (DeepSeek, Google, xAI). Five of the six per-dimension $\alpha$ values on the advisory subset are $\geq 0.85$ (constraint adherence $0.94$, specificity $0.91$, actionability $0.86$, personal fit $0.94$, trade-off awareness $0.92$). The sixth, \emph{risk control}, is $\alpha = 0.78$ on this subset, slightly below the conventional $0.80$ threshold. The drop has a clean explanation: in advisory contexts the no-profile condition triggers generic safety boilerplate that DeepSeek and Grok score as moderate ($2.6$--$2.8$) and Gemini scores as low ($1.9$); the disagreement is concentrated on this single condition rather than spread across the four conditions, and the four-condition ranking remains $\rho = 1.0$ across all judge pairs. We report $\alpha = 0.78$ here rather than substituting the higher full-benchmark figure ($0.82$) because the present paper analyzes the advisory subset alone; readers comparing with the AKM Mother Paper \citep{shao2026akm} will see that the full-benchmark per-dimension $\alpha$ on risk control is slightly higher because fitness and fashion personas pull the no-profile risk-control score upward.

\subsection{The Saturated Rubric}

Both AKM conditions sit near the rubric ceiling ($\approx 4.95/5$ on most dimensions). The current 1--5 scale is not designed to separate \emph{good} structured profiles from \emph{better} structured profiles. The next experiment should adopt forced ranking, per-constraint pass/fail counts, or adversarial personas with deliberately missing information to spread the top end.

\subsection{Why So Much Spread Between AKM and unstructured\_notes?}

The $+8.4$ gap between \textit{akm\_profile} and \textit{unstructured\_notes} on a 30-point scale is large for an LLM-as-judge setting. The most parsimonious reading is that the unstructured note compresses many constraints into ambiguous prose; the structured profile makes each constraint individually addressable. \emph{Constraint adherence} (3.78 vs 4.96) and \emph{trade-off awareness} (2.69 vs 4.96) show the largest gaps, consistent with this reading.

\section{Limitations}
\label{sec:limits}

\paragraph{Synthetic personas.} The simulated user in the elicitation chain has perfect access to the hidden persona JSON. Real users have noisy access to their own state and may distort or omit information. The DaE-elicited number is therefore an upper bound on cold-start recovery, not an estimate of it.

\paragraph{Saturated rubric.} The 1--5 scale leaves only a $0.33$-point gap between the curated and elicited AKM conditions on totals. We cannot reject the null that on a more sensitive instrument the gap is larger.

\paragraph{Single generator.} All conditions share \texttt{deepseek-v4-pro}. Whether the gap persists with weaker or stronger generators is open.

\paragraph{LLM-as-judge.} High inter-judge agreement does not entail human-judgment validity. A small human-rater calibration on a 6--10 persona subset is the natural next step.

\paragraph{Domain coverage.} Five advisory sub-domains do not represent high-stakes decision support (clinical, legal, financial fiduciary). Results should not be transferred to those domains without dedicated evaluation.

\paragraph{Best-answer parsing.} 18 of 90 best-answer selections were not recoverable from one judge's structured output. The dimension scores on which all main results rest are unaffected; only the best-answer count is.

\section{Boundary Conditions and Strategic Scope}

The narrow claim this paper licenses is: \emph{under controlled synthetic personas spanning five advisory sub-domains and three cross-family LLM judges in near-perfect agreement, a single-pass DaE cold-start elicitation produces a structured user profile that matches a hand-curated profile on six downstream advisory-quality dimensions, while both substantially outperform an unstructured first-person note (and trivially outperform no profile at all).}

The paper does not claim production performance on real users, transfer to high-stakes decisions, or universal personalization across cultures.

\section{Conclusion}

The bottleneck in AI advisory is alignment transfer. DaE attacks that bottleneck at the cold-start moment by converting the first interaction into a short structured elicitation pass whose product is a reusable profile. On a 30-persona cross-family three-judge benchmark, the DaE-elicited profile recovers $\geq 98.9\%$ of the hand-curated upper bound on six advisory-quality dimensions, with $\alpha = 0.950$ across judges. The next steps are methodological: ceiling-resistant rubrics, adversarial personas with missing information, and small-scale human calibration. Until those are in place, the empirical claim should be read at the bounded scope above. The mechanism-design claim is more durable: \textbf{convert the cold-start moment into a structured elicitation pass; reuse the resulting profile downstream.}

\paragraph{Code and Data.} Full reproducibility package at \url{https://github.com/sirsws/akm} under \texttt{branches/dae} (mechanism) and \texttt{experiments/synthetic\_profile\_eval/} (benchmark). The advisory subset analyzer is \texttt{subset\_stats.py --tag advisory}.

\bibliographystyle{plainnat}
\bibliography{refs}

\end{document}
